\section{总结与延伸}

\subsection{核心要点回顾}

本讲从最基本的 n-gram 语言模型出发，逐步揭示了现代 LLM 的工作原理和工程实践。核心要点可以归纳为以下几条：

\begin{enumerate}
    \item \textbf{一切始于 next-word prediction}：从 1980 年代的 Bayesian 语言模型到 2025 年的 Gemini，核心任务没有改变——预测下一个 token。所有看似``智能''的行为都是这个简单目标的涌现结果。

    \item \textbf{Scale 带来质变}：参数量（百万 $\rightarrow$ 万亿）和上下文窗口（4 词 $\rightarrow$ 200 万 tokens）的指数级增长，催生了 few-shot learning 等涌现能力。

    \item \textbf{LLM 是训练数据的模糊查询}：模型的输出本质上是训练数据的统计模式重现。Role Prompting 的有效性（``MIT mathematician'' 让数学题正确率提升）印证了这一点——我们只是在条件化输出分布的不同区域。

    \item \textbf{Prompt Engineering 是概率空间操纵}：Role Prompting 条件化分布；Chain-of-Thought 增大错误表面积；格式提示引导模型进入特定的数据模式。

    \item \textbf{Post-training 移动语言模型}：Instruction Tuning、RLHF、Constitutional AI 都是在``合法语言模型''之间移动的手段，本质上是通过更新权重来 shift 概率分布。

    \item \textbf{LLM 有严重的局限性}：幻觉、偏见、Prompt Injection、不守规则、对错误前提不加质疑。使用 LLM 必须配合严格的评估和外部安全机制。

    \item \textbf{Agent = 推理 + 工具使用}：ReAct 融合推理与行动，Toolformer 教模型使用外部工具。这两篇论文奠定了现代 Agent 系统的理论基础。
\end{enumerate}

\subsection{延伸阅读}

\begin{itemize}
    \item \textbf{GPT-3 论文}：Brown et al., ``Language Models are Few-Shot Learners'', NeurIPS 2020 —— 奠定了 few-shot prompting 的基础
    \item \textbf{Chinchilla 论文}：Hoffmann et al., ``Training Compute-Optimal Large Language Models'', 2022 —— 训练数据与参数量的最优比例
    \item \textbf{Chain-of-Thought}：Wei et al., ``Chain-of-Thought Prompting Elicits Reasoning in Large Language Models'', NeurIPS 2022
    \item \textbf{Zero-shot CoT}：Kojima et al., ``Large Language Models are Zero-Shot Reasoners'', NeurIPS 2022 —— ``Let's think step by step'' 的发现
    \item \textbf{LoRA}：Hu et al., ``LoRA: Low-Rank Adaptation of Large Language Models'', ICLR 2022
    \item \textbf{RLHF}：Ouyang et al., ``Training language models to follow instructions with human feedback'', NeurIPS 2022
    \item \textbf{Constitutional AI}：Bai et al., ``Constitutional AI: Harmlessness from AI Feedback'', Anthropic 2022
    \item \textbf{ReAct}：Yao et al., ``ReAct: Synergizing Reasoning and Acting in Language Models'', ICLR 2023
    \item \textbf{Toolformer}：Schick et al., ``Toolformer: Language Models Can Teach Themselves to Use Tools'', NeurIPS 2023
    \item \textbf{LearnPrompting.org}：讲师推荐的 Prompt Engineering 学习资源
\end{itemize}

%% --- Fin del contenido --- %%

\end{document}
